\section{A degree-1081 equation in twelve positive unknowns}
\label{sec:twelve}

\subsection{The polynomial, with no hidden unknowns}
The twelve supplied integers are
\[
 (J,F,\alpha,z,h,i,T,s,w,\eta,\zeta,r)\in\PP^{12}.
\]
Let $(\beta,K_0,M_C,M_F)$ be the fixed tuple of
Section~\ref{sec:coefficients}. In particular, $M_F$ is the shifted mask.
Define the common polynomial expressions
\begin{align*}
 q&=\beta J+1,& X&=wq,&Y&=sq^3,&E&=XY,\\
 k&=\eta+\zeta,&c&=kY+\eta,&a&=E+Y,&H&=4a+3,\\
 \Delta&=a^2+H=(a+2)^2-1,&
 \ell&=EkY,&A&=1+\ell(\ell+k),&G&=1+\Delta c^2.
\end{align*}
The packed rank and two outer expressions are
\begin{align}
 R&=((q-1)(q-F)+\alpha)(q^2-1)+(M_C+qM_F)J,\notag\\
 C_w&=q-F-\alpha,\notag\\
 N_k&=k-hE-R,\qquad
 N_t=(K_0+w)C_w+q-F-z(q-1).\label{eq:outer12}
\end{align}
The normalized strong and auxiliary quantities are
\begin{align*}
 I&=ic^2,&D&=1+\Delta I^2,&Q&=\Delta^2I^2,&d&=Q-1,\\
 v&=cT,&b&=c+RD,&e&=Q(Dv^2+b^2)-1,&B_a&=2Qvb.
\end{align*}
Our replacements for the separating weights are
\begin{equation}
 \boxed{W=T^2+XI^2,\quad L=1+HW,\quad J_\tau=4HW,\quad
 S=X+ac+Hr.}\label{eq:newweights}
\end{equation}
For a polynomial in a formal letter $u$, first reduce it modulo $u^2-g$ to
$p_0+p_1u$, and define
\[
 \Norm_{u^2=g}(p)=p_0^2-gp_1^2.
\]
This is a finite operation on integer polynomials, not a new function in the
Diophantine language. Define
\begin{equation}
 N=\Norm_{u^2=G}\Norm_{\tau^2=A}
 \left(\left[d(S-Lu-J_\tau\tau)^2-H^2e\right]^2
              -H^4B_a^2D\right).
 \label{eq:norm12}
\end{equation}
Then the three presentations are
\begin{align}
 \mathcal S_{12}:&\quad N_k=1,\quad N_t=1,\quad N=0;\label{eq:system12}\\
 P_{12}&=\boxed{N_kN_t(1+N)-1;}\label{eq:product12}\\
 Q_{12}&=(N_k-1)^2+(N_t-1)^2+N^2.\label{eq:sos12}
\end{align}

\subsection{An entirely elementary expansion of the norm}
For readers who prefer a formula containing only addition and multiplication,
the following expressions expand \eqref{eq:norm12}. They also specify a
short shared evaluator. The letter $A_0$ below is unrelated to $A$:
\begin{align*}
 P_0&=QDv^2,&U_0&=Qb^2,&e&=P_0+U_0-1,\\
 V_0&=S^2+L^2G,&V_1&=2SL,&A_0&=J_\tau^2A,\\
 C_0&=d(V_0+A_0)-H^2e,&C_1&=dV_1,\\
 K_f&=4d^2A_0,&Z&=4H^4P_0U_0,\\
 C_{s0}&=C_0^2+GC_1^2,&C_{s1}&=2C_0C_1,\\
 W_0&=C_{s0}-K_fV_0-Z,&W_1&=C_{s1}-K_fV_1,\\
 N_0&=W_0^2+GW_1^2-4ZK_fV_0,\\
 N_1&=2W_0W_1-4ZK_fV_1,\\
 N&=N_0^2-GN_1^2.
\end{align*}
These equations are definitions of abbreviations. They do not add
equations to \eqref{eq:system12} or unknowns to \eqref{eq:product12}.
An exact sparse-polynomial check independently compares this expansion with
the conjugate product; the comparison has 295 monomials in its generic
norm parameters.

\subsection{The inherited arithmetic interface}
We state precisely the reusable facts needed in the new proof. They are
proved in the repository's ordinary blank-input, root-amalgamation and
signed-domain notes, and preserved by the complete Wang coefficient adapter
\cite{p19,p21,arithmetic}.

\begin{proposition}[Weak transport and signed normalized interface]
\label{prop:interface}
With the certified fixed tuple, strictly positive supplied coordinates,
and $N_k=N_t=\epsilon\in\{1,-1\}$, one has
\[
 C_w\ge0,\quad q\ge16,\quad 0<R<q^4-2,\quad E\ge q^4>R+2.
\]
If, in addition, integers $u,\tau,f,y,\gamma$ satisfy
\begin{align*}
 u&=X+ac+H\gamma,& u^2-\Delta c^2&=1,\\
 \tau^2-\ell(\ell+k)&=1,& f^2-\Delta i^2c^4&=1,\\
 Q\bigl(c(Tf-1)-Rf^2\bigr)^2-(Q-1)y^2&=1,
 \end{align*}
then $\epsilon=1$ and a blank K215 halt can be recovered. The recovered
main congruence coordinate and auxiliary roots may initially have signs.
Conversely, a blank halt supplies canonical positive witnesses to the
normalized interface, with all these quantities positive.
\end{proposition}

This is a substantive inherited theorem. In particular, it is not permissible
to cite only its positive-root version after eliminating roots. Its proof
uses signed Pell-rank congruences to identify the main index as $R$, excludes
the negative outer sign and the boundary $C_w=0$, recovers $X=2^R$, and then
decodes the native masks. The ordinary strong-equation formulation follows
from the normalized one by $i_{\mathrm{ord}}=\Delta i$, preserving $Q$.
The full inherited argument, including the signed-quotient issue, is retained
with the archive rather than replaced by finite testing.

\subsection{Coherent roots and a uniform conjugate bound}
The first immediate strengthening of the weak bounds is
\begin{equation}
 k=R+hE+\epsilon\ge R+E-1>2R.\label{eq:k2r}
\end{equation}
Here the positivity of the supplied $h$ is essential. Also
$X\ge16$, $Y\ge4096$, $a=Y(X+1)$, $c>1$, $I=ic^2>c$, and $Q>1$.

Factoring \eqref{eq:norm12} in a splitting field shows that $N=0$ gives
one coherent selection
\begin{equation}
 u^2=G,\quad \tau^2=A,\quad f^2=D,\quad
 dy^2=Q(vf-b)^2-1,\quad S=Lu+J_\tau\tau+Hy.
 \label{eq:coherent12}
\end{equation}
Indeed, the full product is over the eight sign choices for $u,\tau,f$ of
$d(S-Lu-J_\tau\tau)^2-H^2e+H^2B_af$; choose one vanishing factor and divide
its linear expression by $H$ to define $y$. Thus the signs are compatible
by construction. No assertion about independence of square classes is used.

\begin{lemma}[Smaller bound for every auxiliary conjugate]
\label{lem:newbound}
For every complex embedding of a splitting field containing
\eqref{eq:coherent12}, its image $y'$ satisfies
\[
 |y'|<W\sqrt G,\qquad W=T^2+XI^2.
\]
\end{lemma}
\begin{proof}
Every embedding sends $f$ to $\pm\sqrt D$. The auxiliary equation gives
\[
 |y'|<\sqrt{\frac{Q}{Q-1}}\,(v\sqrt D+b).
\]
This also covers nonreal $y'$. The squared expression is real; its absolute
value is bounded by $Q(v\sqrt D+b)^2-1$, which is positive and exceeds one.
The ratio $I>c$ implies the exact inequality
\[
 \frac DG=\frac{1+\Delta I^2}{1+\Delta c^2}<\frac{I^2}{c^2}.
\]
Together with $\sqrt D<(a+2)I$ and $\sqrt G>ac$, this gives
\[
 \frac{v\sqrt D}{\sqrt G}<TI,\qquad
 \frac c{\sqrt G}<\frac1a,\qquad
 \frac D{\sqrt G}<\frac{(a+2)I^2}{c}.
\]
Since $\sqrt{Q/(Q-1)}<3/2$, $R/k<1/2$, and $c>kY$,
\begin{align*}
 \frac{|y'|}{\sqrt G}
 &<\frac32TI+\frac3{2a}+\frac32\frac{R(a+2)I^2}{c}\\
 &<\frac32TI+\frac3{2a}+\frac34(X+2)I^2.
\end{align*}
The last step uses $(a+2)/Y=X+1+2/Y<X+2$. Completing the square gives
$\frac32TI\le T^2+\frac9{16}I^2$. Also $3/(2a)<I^2$. Consequently
\[
 \frac{|y'|}{\sqrt G}
 <T^2+\left(\frac34X+\frac{49}{16}\right)I^2
 <T^2+XI^2,
\]
because $X\ge16$.
\end{proof}

\begin{lemma}[The two root separations]
\label{lem:separations}
For every such conjugate $y'$,
\[
 J_\tau\sqrt A>L\sqrt G+H|y'|,
 \qquad L\sqrt G>H|y'|.
\]
\end{lemma}
\begin{proof}
Since $\sqrt G<(a+2)c$, $\sqrt A>\ell=kXY^2$, and $c<k(Y+1)$,
\[
 \frac{\sqrt G}{\sqrt A}
 <\frac{X+2}{X}\frac{Y+1}{Y}
 \le\frac98\frac{17}{16}=\frac{153}{128}<\frac43.
\]
Even $Y\ge16$ suffices for this deliberately loose estimate. Since $HW>1$,
\[
 4HW\sqrt A>3HW\sqrt G>(1+2HW)\sqrt G
              =L\sqrt G+HW\sqrt G.
\]
Lemma~\ref{lem:newbound} proves the first assertion. The second follows
from $L=1+HW$ and the same bound.
\end{proof}

\subsection{Rational roots become integer roots}
If $\tau$ were irrational, an automorphism of a finite normal splitting
field would send it to $-\tau$. Subtract the image of the last equation in
\eqref{eq:coherent12} from the original. The magnitude of the changed
$\tau$ term is $2J_\tau\sqrt A$, while the other changed terms have total
magnitude at most $2L\sqrt G+2HW\sqrt G$. This contradicts
Lemma~\ref{lem:separations}. Thus $\tau$ is rational and, since its square
is an integer, integral.

If $u$ were irrational, an automorphism sending $u$ to $-u$ would fix the
now rational $\tau$. The same subtraction contradicts
$L\sqrt G>HW\sqrt G$. Hence $u$ is also an integer. The linear relation
in \eqref{eq:coherent12} makes $y$ rational. Moreover,
\[
 f=\frac{e-dy^2}{B_a}
\]
is rational, since $B_a>0$. Its integral square $D$ makes $f$ integral.
It remains to justify integrality of $y$, which does not follow merely
from $Hy\in\ZZ$.

\begin{lemma}[Special auxiliary Pell saturation]
\label{lem:saturation}
If $m\ge2$ and $V$ are integers, $y\in\QQ$, and
\[
 m^2V^2-(m^2-1)y^2=1,
\]
then $y\in\ZZ$.
\end{lemma}
\begin{proof}
Write $y=a/b$ in lowest terms with $b>0$. The equation gives
$b^2\mid m^2-1$. Set $D_0=(m^2-1)/b^2$. It is a positive nonsquare
integer: a square value would make $m^2-1$ a square. We have two positive
integer Pell solutions, after taking absolute values:
\[
 m^2-D_0b^2=1,\qquad (m|V|)^2-D_0|a|^2=1.
\]
Their ordinates cannot vanish. Write the positive powers of the fundamental
integer Pell unit as $\chi_n+\psi_n\sqrt{D_0}$. For some $j,l\ge1$,
\[
 m=\chi_j,\quad b=\psi_j,\quad m|V|=\chi_l,\quad |a|=\psi_l.
\]
The addition formulas imply
$\chi_{n+2j}\equiv-\chi_n\pmod{\chi_j}$ and
$\chi_{2j-n}\equiv-\chi_n\pmod{\chi_j}$. Reducing $l$ modulo $2j$ and
reflecting into $0\le n\le j$, strict increase of $\chi_n$ shows that
$\chi_j\mid\chi_l$ is possible only for $n=j$. Thus $l$ is an odd multiple
of $j$. The Pell multiplication identities then give $\psi_j\mid\psi_l$.
Hence $b\mid a$, and coprimality forces $b=1$.
\end{proof}

Apply Lemma~\ref{lem:saturation} with $m=\Delta I$ and $V=vf-b$. It makes
$y$ integral. Reducing the integral linear relation modulo $H$, and using
$L\equiv1\pmod H$ and $J_\tau\equiv0\pmod H$, gives
$u\equiv X+ac\pmod H$. Thus
$\gamma=(u-X-ac)/H$ is an integer. All the equations of
Proposition~\ref{prop:interface} are recovered, because $f^2=D$ gives
$vf-b=c(Tf-1)-Rf^2$.

\subsection{Soundness of the product and positive completeness}
For any integer $a$, $\Delta=(a+2)^2-1$ is $0$ or $3$ modulo $4$, and
$G=1+\Delta c^2$ is $0$ or $1$ modulo $4$. Therefore
$N=N_0^2-GN_1^2$ cannot be $2$ modulo $4$. At a zero of
$P_{12}=N_kN_t(1+N)-1$, each of the three factors is an integer unit.
The possibility $1+N=-1$ would give $N=-2\equiv2\pmod4$, a contradiction.
Consequently $N=0$ and $N_k=N_t=\epsilon\in\{1,-1\}$.

The preceding reconstruction now applies without having assumed the signs
it eventually proves. Proposition~\ref{prop:interface} forces $\epsilon=1$
and recovers a blank K215 halt. In particular the product has no extra
positive zeros arising from its unit signs.

Conversely, take a canonical positive normalized compiler witness. Let
$u=X+ac+H\gamma$, and take its positive first and auxiliary roots
$\tau,y$. Define
\[
 r=\gamma+Wu+4W\tau+y>0.
\]
Then $S=Lu+4HW\tau+Hy$, so a coherent factor of the norm vanishes.
The outer factors equal one. This constructs a positive zero of the system,
the product, and the sum of squares.

\begin{theorem}[Twelve-variable consistency representations]
The system \eqref{eq:system12}, the polynomial equation $P_{12}=0$, and
the polynomial equation $Q_{12}=0$ have positive integer solutions exactly
when K215 halts on its blank tape. Consequently ZFC proves that the
unsolvability of each is equivalent to $\Con(\ZFC)$.
\end{theorem}

\subsection{Exact degree, circuit size, and the present degree obstruction}
The important exact degrees are
\begin{center}
\begin{tabular}{lrrrrrrrr}
\toprule
Expression & $c$ & $a,H$ & $\Delta$ & $I$ & $D$ & $Q,d$ & $b$ & $e$\\
Degree & 5 & 6 & 12 & 11 & 34 & 46 & 38 & 122\\
\midrule
Expression & $W$ & $L,J_\tau$ & $V_0,A_0$ & $V_1$ & $C_0$ & $C_1$ & $N_0$ & $N$\\
Degree & 24 & 30 & 82 & 41 & 134 & 87 & 536 & 1072\\
\bottomrule
\end{tabular}
\end{center}
The term $H^2e$ has degree $134$, while $d(V_0+A_0)$ has degree $128$.
Thus $C_0$ has unique leading term $-H_{\rm top}^2e_{\rm top}$. The even
coefficient $N_0$ has leading term $C_{0,\rm top}^4$; the odd coefficient
$N_1$ has degree at most $489$, so $GN_1^2$ has degree at most $1000$.
It cannot cancel the degree-$1072$ term of $N_0^2$.

More explicitly, put $q_0=\beta J$ and $k_0=\eta+\zeta$. The norm's
leading homogeneous form is
\begin{equation}
 N_{\rm top}=2^{32}i^{48}k_0^{96}w^{80}s^{176}
                 q_0^{656}(q_0-F)^{16}.\label{eq:leading12}
\end{equation}
It is nonzero for every fixed $\beta>0$. Since
\[
 (N_k)_{\rm top}=-hwsq_0^4,\qquad
 (N_t)_{\rm top}=w(q_0-F-\alpha)-zq_0,
\]
the product has exact degree $1072+7+2=1081$. For example, the coefficient
of
\[
 J^{677}hz\,i^{48}\eta^{96}w^{81}s^{177}
\]
is $2^{32}\beta^{677}\ne0$. The sum of squares has degree $2144$.

The strong block is computed as
$I=ic^2$, $I_2=I^2$, $D_-=\Delta I_2$, $D=D_-+1$, $Q=\Delta D_-$.
It costs the same five gates as the old block, while exposing $I^2$ for
reuse. Similarly $T_2=T^2$, $v^2=c^2T_2$ replaces $v=cT$, $v^2=v\cdot v$
at the same cost; the actual norm circuit only uses $v$ through $v^2$.
This exposes $T^2$ for the weight. These shared computations recover the
extra cost of the new first-root coefficient.
\begin{center}
\begin{tabular}{lrrrr}
\toprule
Presentation & Multiplications & Additions/subtractions & Total & Degree\\
\midrule
Three-equation system & 72 & 40 & 112 & 1072\\
Product & 74 & 42 & 116 & 1081\\
Sum of squares & 75 & 44 & 119 & 2144\\
\bottomrule
\end{tabular}
\end{center}
Once the weights are this small, the dominant degree comes from $H^2e$,
not from the separating weights. Merely reducing those weights further
will not automatically lower this circuit's degree. A further substantial
degree reduction should address the auxiliary constant $b=c+RD$, a
different norm presentation, or a different arithmetic interface.
